The upper half-plane supplies the geometry on which the trace formulas live. Its points may be viewed as complex numbers, as cosets of a compact subgroup of a real matrix group, or as points of a disk. Each description makes a different calculation transparent. The matrix action explains automorphy; the homogeneous-space description separates position from rotation; the metric and its point-pair invariant provide the building blocks for invariant kernels.

We begin with the real group before imposing integrality. This separates facts about hyperbolic geometry from the special arithmetic of the modular group. The first five sections establish the action, Iwasawa coordinates, geodesics, radial integration, and the three types of nonidentity transformations. The last three sections construct the modular fundamental domain, identify its stabilizers and cusp, and introduce Haar measure and its relation to hyperbolic area. The arguments keep the central matrices $\pm I$ visible throughout.

Our geometric quotient uses the effective action of $\operatorname{PSL}_2(\mathbb Z)$ and has hyperbolic area $\pi/3$. A quotient of $\operatorname{SL}_2(\mathbb R)$ by $\operatorname{SL}_2(\mathbb Z)$ has an additional compact direction, whose normalization must be stated before assigning a covolume. We construct Haar measure in Iwasawa coordinates and recover hyperbolic area by integrating over the compact fiber. General uniqueness and left/right properties are proved in \hyperref[v1a-app-a-s2]{Appendix A.2}. We then derive the relation between quotient measures, including the correction needed for subgroups that do not contain $-I$. This will allow later changes of level and viewpoint without silently changing constants.

The chapter proves the elementary geometry and integration identities used here. Spectral theory, automorphic kernel convergence, and trace formulas themselves belong to later chapters. In particular, the invariant kernel construction below is a geometric starting point, not a claim that an arbitrary kernel defines a trace-class operator.

Background references are J. S. Milne's \href{https://www.jmilne.org/math/CourseNotes/MF.pdf}{\emph{Modular Functions and Modular Forms}, Sections 1--2}, and A. V. Sutherland's \href{https://math.mit.edu/classes/18.786/2024/LectureNotes1.pdf}{\emph{Geometry of the complex upper half plane}}. The definitions and calculations below specify this book's conventions independently; in particular, our symbol $\rho$ denotes the left vertex of the standard region.
